\label{sec:conclusion}

We have provided the first systematic analysis of 2030 projected reapportionment
outcomes and their implications for the \textsc{bisect} redistricting platform.
The main findings are:

\begin{enumerate}
\item The projected 2030 Huntington-Hill apportionment produces 8 gaining states
  (TX $+3$, FL $+2$, plus 6 states gaining 1 seat each) and 10 losing states
  (NY $-2$, plus 9 states losing 1 seat each), with 32 states unchanged.
  Texas at $k=41$ and Florida at $k=30$ are the largest redistricting challenges.

\item Projected 2020--2030 county disruption averages 58\%, lower than the observed
  2010--2020 average of 66\%. The lower disruption reflects predictable Sun Belt
  growth patterns and the contraction structure of losing states, both of which
  make optimal county groupings more stable across cycles.

\item Texas $k=41$ triggers the prime-factor fallback in the \texttt{prime-factor}
  structure (since 41 is prime), producing a $20+21$ binary top-level split.
  The ratio-optimal structure (\texttt{ratio-optimal}) is recommended for Texas 2030
  because it evaluates all candidate root splits and selects the GeoSection-optimal
  partition, providing marginal improvements in compactness (PP $+0.02$) and county
  preservation ($+3$ pp) over the standard-bisect fallback.
\end{enumerate}

\subsection{Implications for Q-Series Infrastructure Planning}

The Q.4 findings directly inform the other Q-series tasks:

\begin{itemize}
\item \textbf{Q.3 block-group resolution}: The four newly requiring block-group
  resolution states (CO, NC, and two others based on finalized projections) are
  identified by applying the F.3 rule to Q.4 $k$-values. Q.3 should begin adjacency
  construction planning as soon as Q.4 $k$-values are confirmed.

\item \textbf{Q.2 pre-registration}: The Texas \texttt{ratio-optimal} recommendation
  and the three uncertain states (MT, OR, ID) requiring dual-configuration preparation
  are inputs to Q.2's state-by-state configuration audit.

\item \textbf{Q.1 data update}: States gaining seats (particularly TX, FL, AZ) will
  show the largest demographic changes from 2021 to 2028, making the LODES/ACS update
  most critical for these states.

\subsection{Limitations and Robustness}

The primary limitation is projection uncertainty. Census Bureau medium-series projections
have historically over-estimated growth in states that subsequently experienced economic
slowdowns (California 2000s) and under-estimated growth in fast-growing states
(Arizona, Georgia 2010s). The 2030 projections may miss:

\begin{itemize}
\item Post-pandemic migration reversal from Sun Belt states (if remote work normalization
  reverses the 2020--2022 domestic migration surge).
\item Climate-driven migration from vulnerable coastal areas (FL, TX Gulf Coast).
\item Immigration policy effects on Hispanic growth rates (major driver of TX, AZ, FL growth).
\end{itemize}

We recommend treating the Q.4 seat allocations as planning inputs with uncertainty
ranges (Table~\ref{tab:uncertain}) rather than point predictions. Infrastructure
preparation should proceed for the median projection while maintaining flexibility
for adjacent seat allocations in the three boundary states.

\bigskip
\noindent\textit{The author thanks the Census Bureau's Population Division for public
access to the 2023 National Population Projections and the Defense Manpower Data Center
for 2022 state-of-legal-residence counts used for the overseas adjustment.
All errors are the author's own.}
